\section{Practical Implications and Deployment Considerations}
\label{sec:praxis}

\subsection{Prompt Injection as a System-Level Risk}

A key conclusion from the research surveyed in this paper is that prompt
injection should not be treated as an isolated model weakness. In practical
deployments, the LLM is only one component within a larger system that may
include retrieval pipelines, external APIs, memory systems, and autonomous
tool execution. Consequently, a successful injection does not stop at the
model itself, but propagates into whatever capabilities the surrounding system
exposes.

This changes the nature of the threat. The relevant question is not simply
whether the model can be manipulated, but what actions become possible once
manipulation occurs. As AgentDojo demonstrates, increasingly autonomous
systems also become increasingly attractive targets
~\parencite{debenedetti2024agentdojo}. Real-world incidents, including attacks
against Slack AI, illustrate that prompt injection can lead to unintended
information disclosure and actions extending beyond the model
~\parencite{chen2025secalign}.

\subsection{Defence in Depth}

Since no single defence mechanism eliminates the attack surface, current
research increasingly favours a layered approach. The goal is not to prevent
every successful injection, but to ensure that a successful attack cannot
easily escalate into a critical incident.

A practical deployment may combine model-level defences with restricted tool
permissions, sandboxed execution environments, monitoring, and human approval
for sensitive operations. No individual layer is sufficient, but together
they reduce the probability that an attacker can progress from injection to
impact.

Human oversight remains particularly important for high-risk tasks. Requiring
user confirmation before irreversible actions and limiting tool access help
reduce the consequences of successful attacks. These principles mirror the
established defence-in-depth philosophy long used in traditional software
security~\parencite{chen2025secalign}.

\subsection{Deployment Trade-Offs}

Secure deployment inevitably involves trade-offs between security, usability,
and operational complexity. Architectural approaches based on information flow
control provide stronger guarantees, but often require substantial changes to
existing systems~\parencite{wu2024systemlevel, zhong2025rtbas}. Fine-tuning
approaches such as StruQ and SecAlign are easier to integrate, but their
robustness remains limited against adaptive attacks
~\parencite{yang2025checkpointgcg}.

Detection systems are operationally lightweight, but introduce false positives
and can themselves be bypassed~\parencite{jia2025critical}. Consequently, no
currently available approach offers an ideal combination of security,
flexibility, and ease of deployment.

Commercial environments further complicate this balance. Pressures for rapid
deployment and ease of use often conflict with conservative security design.
As LLM systems gain more autonomy and access to sensitive operations, the cost
of underestimating prompt injection risks is likely to increase.

\subsection{Future Directions}

Several directions emerge from the limitations identified throughout this
paper. One promising avenue is stronger architectural isolation, shifting away
from relying solely on model behaviour and instead enforcing trust boundaries
through system design~\parencite{wu2024systemlevel,
debenedetti2025defeating}.

At the same time, evaluation methodologies must evolve to treat adaptive
attackers as the baseline rather than the exception. More realistic benchmarks
and standardised utility metrics are necessary if robustness claims are to
reflect practical security~\parencite{debenedetti2024agentdojo,
fspib2025, jia2025critical}.

Finally, increasing model capability does not automatically imply improved
security. More capable systems may also be more effective at carrying out
malicious instructions, suggesting that capability and robustness do not
necessarily scale together~\parencite{debenedetti2024agentdojo}.

Ultimately, prompt injection cannot currently be considered a solved problem.
While practical safeguards can significantly reduce risk, they do not
eliminate it entirely. This raises the broader question of how effective
existing defence strategies truly are and where their fundamental limitations
lie, which is addressed in the concluding section.
